\begin{textsample}{NEG Dim 8 – global\_north\_2024\_11 – Score -1.00 – global\_north\_2024\_11/118142377.txt}  \label{ex:f8_neg_392}
What Israel and Hezbollah ' s ceasefire deal will look like in action In the end, Israel and Hezbollah concluded that they could gain more through negotiations than they could on the battlefield.
Nov. 26, 2024, 8:55 PM UTC By Daniel R.
DePetris, fellow at Defense Priorities President Joe Biden ' s announcement on Tuesday of a ceasefire agreement between Israel and Hezbollah couldn't have come soon enough for Lebanon, a country in the midst of a yearslong economic crisis and intense political paralysis.
The war, which began on Oct. 8, 2023, as a series of hostile exchanges across the Israel-Lebanon border and escalated into a heavy Israeli air and ground campaign in Southern Lebanon, has destroyed entire communities and turned some of Beirut ' s districts into a war zone.
Hours before the U.S.-brokered deal with Lebanon was announced, Israel struck multiple Hezbollah targets in Beirut in what was no doubt a message to the Lebanese militia : Israel can sustain the conflict for as long as it sees fit.
Hours @ @ @ @ @ @ @ @ @ @ Hezbollah targets in Beirut in what was no doubt a message to the Lebanese militia : Israel can sustain the conflict for as long as it sees fit.
In the end, Israel and Hezbollah concluded that they could gain more through negotiations than they could on the battlefield.
The agreement is a recitation of United Nations Security Council Resolution 1701, which ended a previous monthlong war between the two adversaries more than 17 years ago but was viewed by all the parties involved, Israel in particular, as a lackluster initiative that wasn't enforced.
The current deal seeks to strengthen UNSCR 1701 by adding stronger monitoring.
During a 60-day ceasefire, Israeli troops will withdraw from Southern Lebanon, Hezbollah will do the same, and the Lebanese army will re-deploy to the area.
Meanwhile, the approximately 60,000 Israelis who have been displaced in northern Israel will get to return home.
In effect, the deal allows both Israel and Hezbollah to claim victory ; Israeli Prime Minister Benjamin Netanyahu can boast that Hezbollah ' s military capacity has @ @ @ @ @ @ @ @ @ @ drove Israeli forces out of Lebanon.
Over the long term, the pause in fighting is designed to give Israel and Lebanon the time and space to officially demarcate their shared border.
Yet, at the risk of sounding like a pessimist, just getting to that point would be an achievement.
A lot can go wrong between now and then.
After all, signing an agreement means nothing if it isn't implemented.
There are any number of ways the agreement can go sideways.
First and foremost, the question of whether Hezbollah will actually withdraw north of the Litani River, approximately 20 miles from the Israeli-Lebanese border, is very much in question.
Southern Lebanon is Hezbollah ' s support base ; the militia is a core part of the social fabric in the region, its fighters have homes and families there, and the small towns and villages dotting the area have often been given the short end of the stick from the Lebanese government, which has proven incapable of delivering social services or even @ @ @ @ @ @ @ @ @ @ weapons caches further north, but the idea that tens of thousands of Hezbollah fighters will uproot their lives is difficult to believe.
In this case, Israel will then be forced with a choice : renew military operations and risk the resumption of war, or loosen enforcement and risk Hezbollah maintaining its power base.
Second, is the Lebanese army capable of patrolling Southern Lebanon to Israel ' s satisfaction?
While the Lebanese army is a well-respected institution inside the country and crosses the usual sectarian divisions that have defined Lebanese political life for decades, it ' s also arguably the weakest military in the Middle East.
At 70,000-80,000 soldiers, the Lebanese army is smaller than Hezbollah.
The roughly \$3 billion in U.S. military aid since 2006 has barely kept the Lebanese army afloat.
The defense systems you would expect a modern military to possess -- air defenses, fighter and bomber aircraft, patrol vessels, various air-to-air and air-to-ground munitions -- are either in short supplyor nonexistent.
Due in large part to Lebanon ' s financial crisis, some @ @ @ @ @ @ @ @ @ @ support themselves and their families.
Israel understands all this but is nevertheless unlikely to be very patient.
If the Lebanese army is unable or unwilling to do the job of clearing Hezbollah forces out and confronting the remnants that remain, the Israeli army will do it for them as Netanyahu made abundantly clear during his remarks upon announcing the ceasefire.
This, in effect, would negate the ceasefire and risk plunging the country into war again.
Assuming the ceasefire in Lebanon sticks, Israeli and U.S. officials are hopeful it will change Hamas ' calculations about continuing its war with Israel in Gaza."What Hamas wanted was support from Hezbollah and others,"an Israeli official told the Times of Israel."Once you cut the connection, you have the ability to reach a deal.
It ' s a strategic achievement.
Hamas is alone."But this sounds more like wishful thinking than reality.
Hamas has experienced the most destructive war with Israel in its 37-year history, with tens of thousands of its @ @ @ @ @ @ @ @ @ @ control in Gaza at its weakest since it kicked the Palestinian Authority out of the coastal territory in 2007.
Even so, it ' s bottom-line negotiating position remains unchanged : If Israel wants to retrieve the rest of its hostages, it must withdraw entirely from Gaza and end the war permanently.
Hamas ' strategy doesn't depend on Hezbollah, so the notion it will adopt Hezbollah ' s position now that it is out of the fight is fanciful at best.
If all goes according to plan, Lebanon will now have a chance to rebuild.
But how long the peace will stick is another matter entirely.
Daniel R.
DePetris Daniel R.
DePetris is a fellow at Defense Priorities and a \textbf{syndicated} foreign affairs columnist at the Chicago Tribune.

% matched lemmas: syndicated
\end{textsample}
